\documentclass[10pt,a4paper]{amsart}
\usepackage[margin=19mm]{geometry}
\usepackage{amsmath,amssymb,mathtools}
\usepackage[scr=rsfs,cal=euler]{mathalfa}
\usepackage{xfrac}
\usepackage[unicode,hidelinks,bookmarks=false]{hyperref}
\newcommand{\ie}{i.e.\ }
\newcommand{\cf}{cf.\ }
\newcommand{\resp}{respectively\ }
\newcommand{\Subsection}{\subsection}
\providecommand{\nobreakdash}{\nobreak\hbox{-}}
\setlength{\emergencystretch}{2em}
\multlinegap0pt
\usepackage{bm}
\usepackage{bbm}
\usepackage{dsfont}
\usepackage{xspace}
\usepackage{cancel}
\usepackage{mathtools}
\usepackage{cleveref}
\usepackage{amsthm}
\makeatletter
\@ifundefined{theorem}{\newtheorem{theorem}{Theorem}}{}
\@ifundefined{lemma}{\newtheorem{lemma}[theorem]{Lemma}}{}
\@ifundefined{proposition}{\newtheorem{proposition}[theorem]{Proposition}}{}
\makeatother
\makeatletter
\DeclareMathOperator{\Con}{\mathrm{Con}}
\DeclareMathOperator{\Inc}{\mathrm{Inc}}
\DeclareMathOperator{\Pro}{\mathrm{Pro}}
\newcommand{\as}{\mathcal{A}}
\expandafter\ifx\csname customthm\endcsname\relax
\newenvironment{customthm}[1]
{\renewcommand\theinnercustomthm{#1}\innercustomthm}{\endinnercustomthm}
\fi
\newcommand{\tot}{\mathrm{Tot}}
\makeatother
\title[Orbitmesy and promotion on self-dual posets]{Orbitmesy and promotion on self-dual posets}
\author{Esther Banaian, Emily Barnard, Sunita Chepuri, Jessica Striker}
\date{Source: arXiv:2508.19440v1 (CC BY 4.0)}
\begin{document}
\maketitle

\noindent\emph{Source and licence.} Orbitmesy and promotion on self-dual posets. Version arXiv:2508.19440v1 (CC BY 4.0), distributed under the Creative Commons Attribution 4.0 International licence, as recorded in the arXiv licence metadata for this exact version and shown on the abstract page https://arxiv.org/abs/2508.19440v1. Full rights notice: licenses/arxiv-2508.19440-CC-BY-4.0.txt.\par\smallskip

\noindent\textit{Editorial scope.} Counted unit: the Proposition whose source label is \texttt{prop:Stable} in the article (source0.tex lines 1675--1677, source label \texttt{prop:Stable}), delivered together with the article's own complete original proof of it (source0.tex lines 1679--1726, one \texttt{proof} environment of 48 lines). No mathematics is written, repaired or re-derived: statement and proof are the article's own text line for line. Required context retained uncounted: prop:ProRestricts (lines 821--823); thm:OrderComparedToDeflation (lines 973--978); lem:GlobalAverage (lines 1577--1587); linear\_ext (lines 1595--1596). Every definition, display and label of those lines is retained unchanged. Omitted from this package and not redistributed: 1--820, 824--972, 979--1576, 1588--1594, 1597--1674, 1678--1678, 1727--2096. Nothing inside a retained excerpt is omitted or altered: every retained body is delivered exactly as the article's lines are written, so no omission receipt of any kind accompanies this package. Citations: the retained excerpts carry no citation command at all, so no bibliography entry of the article is transcribed and no reference is redistributed. Reference adaptation: none was needed and none was made; every reference inside the delivered unit resolves inside it, so no unresolved reference or citation is produced. Printed slips kept literal: no printed word, symbol, hypothesis or number of the counted statement or of its proof was altered. Layout and class: the article's own class is \texttt{amsart} with packages including geometry, subfigure, verbatim, amsmath, amssymb, amsfonts, amsthm, hyperref, mathtools, tikz, enumitem, makecell, mathdots, cleveref; the wrapper is the lane's pinned \texttt{amsart} editorial wrapper at 10pt on A4, which restates the theorem-like environments the retained text needs and adds the article's own macro definitions that the retained text needs (\texttt{\textbackslash Con}, \texttt{\textbackslash Inc}, \texttt{\textbackslash Pro}, \texttt{\textbackslash as}, \texttt{\textbackslash tot}), with extra wrapper packages (bm, bbm, dsfont, xspace, cancel, mathtools, cleveref). The delivered numbering is the unit's own. Assets: no retained span contains a figure environment, a graphics-inclusion command, a PDF-page inclusion, a table or a TikZ picture; no image file is copied into this package, the package declares no asset dependency and no image file is redistributed with it. Licence: version arXiv:2508.19440v1 of this article is distributed under the Creative Commons Attribution 4.0 International licence, as recorded in the arXiv licence metadata for this exact version and shown on the abstract page https://arxiv.org/abs/2508.19440v1. A full rights notice is delivered at licenses/arxiv-2508.19440-CC-BY-4.0.txt. Characters: every retained excerpt is pure ASCII, so the delivered engine, XeLaTeX with its default Latin Modern text font, reports no missing character.
\begin{proposition}\label{prop:ProRestricts}
Let $f \in \Inc^q(P)$ and $\overline{f}$ be $r$-packed. If the first entry in $\Con(f)$ is~$1$, we have $\overline{\Pro(f)} = \Pro(\overline{f})$. Otherwise, we have $\overline{\Pro(f)} = \overline{f}$.
\end{proposition}

\begin{theorem}\label{thm:OrderComparedToDeflation}
Let $f \in \Inc^q(P)$ such that $\Con(f)$ has period $\ell$. Suppose $\overline{f}$ is $r$-packed and  
the size of the promotion orbit containing $\overline{f}$ is $\tau$. Then, the promotion orbit of $f$ has size \[
\frac{\tau \ell}{\gcd(r\ell/q,\tau)}.
\]
\end{theorem}

\begin{lemma}\label{lem:GlobalAverage}
Let $P$ be a self-dual poset and consider all elements of $\Inc^q(P)$. 
\begin{enumerate}
    \item For all $x \in P$, the global average of the antipodal sum statistic with respect to $x$ is $q+1$. That is, \[
\frac{1}{\vert P\vert}\sum_{f \in \Inc^q(P)} \as_x(f) = q+1.
    \]
    \item The global average of the total sum statistic is $\frac{\vert P \vert}{2} (q+1)$. That is, \[
\frac{1}{\vert P \vert } \sum_{f \in \Inc^q(P)} \tot(f) = \frac{\vert P \vert}{2}(q+1).
    \]
\end{enumerate}
\end{lemma}

\begin{proposition}\label{linear_ext} Given a self-dual poset $P$, let $f \in \Inc^q(P)$ such that $\overline{f}$ is a linear extension. Moreover, given $r,\ell,\tau$ as in \Cref{thm:OrderComparedToDeflation}, suppose $\gcd(r\ell/q,\tau) = 1$. If $\mathcal{O}$ is the promotion orbit of $\Inc^q(P)$ containing $f$, then the orbit $\mathcal{O}$ is orbitmesic with respect to the total sum statistic $\tot$.
\end{proposition}

\begin{proposition}\label{prop:Stable}
Let $P$ be a self-dual poset and let $x \in P$. Let $f \in \Inc^q(P)$ be such that $r,\ell,\tau$ from \Cref{thm:OrderComparedToDeflation} satisfy $\gcd(r\ell/q,\tau) = 1$. Let $\mathcal{O}$ be the promotion orbits of elements of $\Inc^q(P)$ containing $f$. If the deflated orbit $\overline{\mathcal{O}}$ is $x$-stable, then $\mathcal{O}$ exhibits orbitmesy with respect to the antipodal sum statistic $\mathcal{A}_x$. 
\end{proposition}

\begin{proof}
As in the proof of \Cref{linear_ext}, the condition on the gcd implies that  $\mathcal{O}$ contains every possible pair of a labeling in $\overline{\mathcal{O}}$ and a cyclic rotation of $\Con(f)$. 
We begin by assuming that $\ell$, the period of $\Con(f)$, is equal to $q$.
We can use the description of $\mathcal{O}$ as pairs of content words and packed labelings to break up the computation of $\mathcal{A}_x$ for all labelings in $\mathcal{O}$. 

We assume we have chosen $f$ so that $\Con(f) = (c_1,\ldots,c_q)$ has $c_q = 1$. Define $s_1,\ldots,s_r$ to be the nonnegative integers such that the $r$ values of $1$ in $\Con(f)$ occur at $s_1+1,s_1+s_2 + 2,\ldots,s_1 + \cdots + s_{r-1} + s_r + r$. That is, $\Con(f)$ is $s_1$ 0's followed by a 1, then $s_2$ 0's followed by a 1, and so on.

Let $j$ and $k$ be such that $\overline{f}(x) = j$ and $\overline{f}(\kappa(x)) = k$. With our notation, this means $f(x) = (\sum_{i=1}^j s_i) + j$ and $f(\kappa(x)) = (\sum_{i=1}^k s_i) + k$. For all $t \leq s_1$ (equivalently, for all $t$ such that $c_1 = c_2 = \cdots = c_t = 0$), from \Cref{prop:ProRestricts} we know $\Pro^t(f)(x) = (\sum_{i=1}^j s_i) + j - t$ and $\Pro^t(f)(\kappa(x)) = (\sum_{i=1}^k s_i) + k - t$.  Therefore, we have
\begin{align*}
\sum_{i=0}^{s_1} \mathcal{A}_x(\Pro^i(f)) &= (s_1 + 1)(\sum_{i=1}^{j} s_i + \sum_{i=1}^{k} s_i + j + k) -2 - 4 - \cdots - 2s_1\\
&= (s_1 + 1)(\sum_{i=1}^{j} s_i + \sum_{i=1}^{k} s_i + j + k) -s_1(s_1+1)\\
& = (s_1 + 1)(\sum_{i=1}^{j} s_i + \sum_{i=1}^{k} s_i + j + k-s_1)\\
& = (s_1 + 1)(s_1 + \sum_{i=2}^{j} s_i + \sum_{i=2}^{k} s_i + j + k).\\
\end{align*}

The computation is similar for all other possible deflations. Let $f_2,\ldots,f_\tau$ be the $\tau-1$ other labelings in $\mathcal{O}$ with content word equal to $\Con(f)$. Let $j_z,k_z$ be such that $f_z(x) = j_z$ and $f_z(\kappa(x)) = k_z$. Let $f_1$ be our original labeling $f$ so that $j_1 = j$ and $k_1 = k$ from the previous summation. Then, we compute
\begin{align}
\sum_{z = 1}^\tau \sum_{i=0}^{s_1} \mathcal{A}_x(\Pro^i(f_z)) = (s_1 + 1)\big(a_1 s_1 + a_2 s_2 + \cdots + a_rs_r + \sum_{z=1}^\tau j_z + k_z)\big)\label{eq:SumFirstSection}
\end{align}
where $a_1 = \tau$ and for $i > 1$, $a_i$ is equal to the number of values in the multiset $\{j_z,k_z\}_{z=1}^r$ which are at least $i$. 
Note that the stable condition implies $\sum_{z=1}^\tau (j_z + k_z) = \tau(r+1)$.  This means \begin{align}
 \sum_{i=1}^r a_i=\tau+\sum_{i=2}^r a_i=\tau+\sum_{z=1}^\tau ((j_z-1)+(k_z-1))=\sum_{z=1}^\tau (j_z+k_z)-\tau= \tau (r+1) - \tau = \tau r.  \label{eq:SimplifySumai} 
\end{align}


Let $g$ be the labeling which deflates to $\overline{f}$ and has content word $\rho^{s_1+1}(\Con(f))$. As noted in the first sentence of this proof, we know $g$ is in the orbit $\mathcal{O}$.  By the definition of the numbers $s_i$,  we know that $\mathcal{A}_x(g) = (\sum_{i=2}^{j+1} s_i) + j + (\sum_{i=2}^{k+1} s_i) + k$.

We can compute $\sum_{i=s_1 + 1}^{s_1 + s_2 + 1} \mathcal{A}_x(\Pro^i(g))$ analogously to the computation of $\sum_{i=0}^{s_1} \mathcal{A}_x(\Pro^i(f))$, but with the values $s_1,\ldots,s_r$ cyclically rotated. Therefore, $\sum_{h \in \mathcal{O}} \mathcal{A}_x(h)$ is the result of summing the $r$ terms given by all cyclic shifts of \Cref{eq:SumFirstSection}. Explicitly, this is \begin{align*}
\sum_{h \in \mathcal{O}}\mathcal{A}_x(h) &= (s_1 + 1)(\tau s_1 + a_2 s_2 + \cdots + a_{r-1}s_{r-1} + a_rs_r + \tau (r+1))\\
& + (s_2 + 1)(\tau s_2 + a_2s_3 + \cdots + a_{r-1}s_r + a_r s_1 + \tau (r+1))\\
& \cdots \\
& +(s_r + 1)(\tau s_r + a_2 s_1 + \cdots + a_{r-1}s_{r-2} + a_rs_{r-1} + \tau (r+1)).
\end{align*}

Treating this as a polynomial in the variables $s_1,\ldots,s_r$ and analyzing the coefficients, we have \begin{align*}
\sum_{h \in \mathcal{O}} \mathcal{A}_x(h) &= \tau(s_1^2 + \cdots + s_r^2) + (a_2+a_r)(s_1s_2 + s_2s_3 + \cdots + s_rs_1) \\
&+ (a_3 + a_{r-1})(s_1s_3 + s_2s_4 + \cdots + s_{r-1}s_1) + \cdots
+ \tau(2r+1) (s_1 + \cdots + s_r) + \tau r (r+1),
\end{align*}
where we use \Cref{eq:SimplifySumai} to analyze for the coefficient of linear terms. As $a_r$ is the number of values in $\{j_z,k_z\}_z$ which are equal to $r$, we have $a_r=2\tau - \# \{j_z,k_z\}_z \cap [1,r-1]$ and by the stable condition this is also $2\tau - \# \{j_z,k_z\}_z \cap [2,r]$. This latter quantity is $2\tau - a_2$ because $\# \{j_z,k_z\}_z \cap [2,r]$ is the number of elements in $\{j_z,k_z\}_z$ that are at least 2. This means $a_2 + a_r = 2\tau$. We can proceed similarly for the other coefficients of the squarefree quadratic terms. Therefore, we have \begin{align*}
\sum_{h\in\mathcal{O}}\mathcal{A}_x(h) &= \tau(s_1^2 + \cdots + s_r^2) + 2\tau(s_1s_2 + s_1s_3 + \cdots + s_{r-1}s_r) + (\tau(2r+1) )(s_1 + \cdots + s_r) \\&+ \tau r (r+1)\\
&= \tau (s_1 + s_2 + \cdots + s_r + r)(s_1 + s_2 + \cdots + s_r + r+1)\\
&= \tau q (q+1).
\end{align*}
Since the size of $\mathcal{O}$ is $\tau q$ by Theorem \ref{thm:OrderComparedToDeflation}, the average value of $\mathcal{A}_x$ on $\mathcal{O}$ is $q+1$. By \Cref{lem:GlobalAverage} this is equal to the global average of $\mathcal{A}_x$ on $\Inc^q(P)$. 

Now, if $\ell$ is a proper divisor of $q$, this computation counts each $h\in\mathcal{O}$ exactly $\frac{q}{\ell}$ times. In this case, the size of $\mathcal{O}$ is $\tau \ell$, and the result again follows.
\end{proof}

\end{document}
